
%\needspace{4\baselineskip}
\color{uniblue}{\subsection{\raggedright{Контроллер присоединения (КП)}}}
\color{black}

\begin{enumerate}[label=\arabic{section}.\arabic{subsection}.\arabic{enumi}, labelsep=4pt, leftmargin=0pt, itemindent=57pt, itemsep=0pt, parsep=5pt] % если что можно крутить itemsep

%%%%%%%%%%%%%%%%%%%%%%%%%%%%%%%%%%%%%%%%%%%%%%%%%%%%%%%%%%%%%%%%%%%%%%  ОБЩИЕ СВЕДЕНИЯ %%%%%%%%%%%%%%%%%%%%%%%%%%%%%%%%%%%%%%%%%%%%%%%%%%
\item Общие сведения

\begin{enumerate}[label=\arabic{section}.\arabic{subsection}.\arabic{enumi}.\arabic*, labelsep=4pt, leftmargin=0em, itemindent=65pt, parsep=0pt]

\item
КП предназначен:
\begin{itemize}
\item для оперативного управления коммутационными аппаратами; 
\item для формирования сигнала положения коммутационного аппарата для целей АСУ ТП; 
\item для реализации оперативной блокировки управляемых коммутационных аппаратов.
\end{itemize}

\item
В устройстве контроллер присоединения представлен функциональным блоком «КП». В функциональном блоке заложена возможность оперативного управления следующими коммутационными аппаратами:
\begin{itemize}
	\item заземляющий нож;
	\item выкатной элемент.
\end{itemize}

\item
При выполнении коммутационной операции ЗН или ВЭ информация об этом процессе будет отображаться общим выходным сигналом «КП: Идет переключение». Функциональная схема формирования общего сигнала представлена на рисунке \ref{kp:img_perkl}.

\vspace{3mm}
\begin{figure}[H]
\centering
\includegraphics[width=0.6\textwidth,height=0.6\textheight,keepaspectratio]{\fbpath/214. КП ТН/_latex/img/2.pdf}
\captionof{figure}{Функциональная схема формирования общего сигнала переключения коммутационного аппарата}\label{kp:img_perkl}
\end{figure}

\item
Ввод функционального блока в работу на этапе параметрирования устройства осуществляется программным переключателем SGF1 «Ввод\_функции». Введенное состояние функционального блока соответствует активному состоянию сигнала «КП: Ввод».
\item
Функциональный блок может быть выведен из работы оперативно путем активации сигнала «ОВ КП», о чем сигнализирует активное состояние сигнала «КП: Оперативный вывод».
\item
Логика управления режимами работы функционального блока выполнена в соответствии с рисунком \ref{kp:img_gen}. В таблице \ref{kp:tbl_gen} приведены параметры, необходимые для настройки функционального блока. 

\vspace{3mm}
\begin{figure}[H]
\centering
\includegraphics[width=1\textwidth,height=1\textheight,keepaspectratio]{\fbpath/214. КП ТН/_latex/img/1.pdf}
\captionof{figure}{Функциональная схема управления режимами работы КП}\label{kp:img_gen}
\end{figure}

\small
\begin{longtable}{|>{\centering\arraybackslash}m{5.3cm}|>{\centering\arraybackslash}m{3.3cm}|>{\centering\arraybackslash}m{4.2cm}|>{\centering\arraybackslash}m{1.8cm}|>{\centering\arraybackslash}m{1cm}|}
\caption{Параметры для настройки КП\hfill\vspace{-0.5\baselineskip}}\label{kp:tbl_gen}\\ 
\hline
\rowcolor{gray!30}
Параметр (Параметр на ИЧМ) & Условное обозначение на схеме & Значение/ Диапазон & Единица измерения & Шаг \\ 
\hline
\endfirsthead
\caption*{\hspace{3pt}\emph{Продолжение таблицы \ref{kp:tbl_gen}\hfill\vspace{-0.5\baselineskip}}} \\ % сделано по ГОСТ 2.105 п.6.8.7
\hline
\rowcolor{gray!30}
Параметр (Параметр на ИЧМ) & Условное обозначение на схеме & Значение/ Диапазон & Единица измерения & Шаг \\ 
%\hline
\endhead
%\hline
\endfoot
\hline
\endlastfoot
%==+t1*LLN0|VT_SWCTRL
\multicolumn{5}{|c|}{ УВЭ } \\ \hline 
\centering Ввод функции в работу (Ввод\_функции)  & \centering SGF1 & \centering Не предусмотрено \\ Предусмотрено & \centering -- & \centering \arraybackslash -- \\
\hline
\centering Допустимое время переключения (Tперекл)  & \centering T1 & \centering 50...100000 & \centering мс & \centering \arraybackslash 10 \\
\hline
\centering Выдержка времени подавления выдачи положения "Не определено" (Тблк)  & \centering T2 & \centering 50...100000 & \centering мс & \centering \arraybackslash 10 \\
\hline
\multicolumn{5}{|c|}{ УЗН } \\ \hline 
\centering Ввод функции в работу (Ввод\_функции)  & \centering SGF1 & \centering Не предусмотрено \\ Предусмотрено & \centering -- & \centering \arraybackslash -- \\
\hline
\centering Допустимое время переключения (Tперекл)  & \centering T1 & \centering 50...100000 & \centering мс & \centering \arraybackslash 10 \\
\hline
\centering Выдержка времени подавления выдачи положения "Не определено" (Тблк)  & \centering T2 & \centering 50...100000 & \centering мс & \centering \arraybackslash 10 \\
\hline
\multicolumn{5}{|c|}{ ОБР ЗН } \\ \hline 
\centering Контроль KA1 (Контр\_KA1)  & \centering SGF1 & \centering Не предусмотрено \\ Предусмотрено & \centering -- & \centering \arraybackslash -- \\
\hline
\centering Контроль KA2 (Контр\_KA2)  & \centering SGF2 & \centering Не предусмотрено \\ Предусмотрено & \centering -- & \centering \arraybackslash -- \\
\hline
\centering Контроль KA3 (Контр\_KA3)  & \centering SGF3 & \centering Не предусмотрено \\ Предусмотрено & \centering -- & \centering \arraybackslash -- \\
\hline
\centering Контроль KA4 (Контр\_KA4)  & \centering SGF4 & \centering Не предусмотрено \\ Предусмотрено & \centering -- & \centering \arraybackslash -- \\
\hline
\centering Контроль KA5 (Контр\_KA5)  & \centering SGF5 & \centering Не предусмотрено \\ Предусмотрено & \centering -- & \centering \arraybackslash -- \\
\hline
\centering Контроль KA6 (Контр\_KA6)  & \centering SGF6 & \centering Не предусмотрено \\ Предусмотрено & \centering -- & \centering \arraybackslash -- \\
\hline
\centering Контроль KA7 (Контр\_KA7)  & \centering SGF7 & \centering Не предусмотрено \\ Предусмотрено & \centering -- & \centering \arraybackslash -- \\
\hline
\centering Контроль KA8 (Контр\_KA8)  & \centering SGF8 & \centering Не предусмотрено \\ Предусмотрено & \centering -- & \centering \arraybackslash -- \\
\hline
\centering Контроль KA9 (Контр\_KA9)  & \centering SGF9 & \centering Не предусмотрено \\ Предусмотрено & \centering -- & \centering \arraybackslash -- \\
\hline
\centering Контроль KA10 (Контр\_KA10)  & \centering SGF10 & \centering Не предусмотрено \\ Предусмотрено & \centering -- & \centering \arraybackslash -- \\
\hline
\centering Контроль KA11 (Контр\_KA11)  & \centering SGF11 & \centering Не предусмотрено \\ Предусмотрено & \centering -- & \centering \arraybackslash -- \\
\hline
\centering Контроль KA12 (Контр\_KA12)  & \centering SGF12 & \centering Не предусмотрено \\ Предусмотрено & \centering -- & \centering \arraybackslash -- \\
\hline
\centering Контроль KA13 (Контр\_KA13)  & \centering SGF13 & \centering Не предусмотрено \\ Предусмотрено & \centering -- & \centering \arraybackslash -- \\
\hline
\centering Контроль KA14 (Контр\_KA14)  & \centering SGF14 & \centering Не предусмотрено \\ Предусмотрено & \centering -- & \centering \arraybackslash -- \\
\hline
\centering Контроль KA15 (Контр\_KA15)  & \centering SGF15 & \centering Не предусмотрено \\ Предусмотрено & \centering -- & \centering \arraybackslash -- \\
\hline
\centering Контроль KA16 (Контр\_KA16)  & \centering SGF16 & \centering Не предусмотрено \\ Предусмотрено & \centering -- & \centering \arraybackslash -- \\
\hline
\centering Контроль KA17 (Контр\_KA17)  & \centering SGF17 & \centering Не предусмотрено \\ Предусмотрено & \centering -- & \centering \arraybackslash -- \\
\hline
\centering Контроль KA18 (Контр\_KA18)  & \centering SGF18 & \centering Не предусмотрено \\ Предусмотрено & \centering -- & \centering \arraybackslash -- \\
\hline
\centering Контроль KA19 (Контр\_KA19)  & \centering SGF19 & \centering Не предусмотрено \\ Предусмотрено & \centering -- & \centering \arraybackslash -- \\
\hline
\centering Контроль KA20 (Контр\_KA20)  & \centering SGF20 & \centering Не предусмотрено \\ Предусмотрено & \centering -- & \centering \arraybackslash -- \\
\hline
\centering Контроль KA21 (Контр\_KA21)  & \centering SGF21 & \centering Не предусмотрено \\ Предусмотрено & \centering -- & \centering \arraybackslash -- \\
\hline
\centering Контроль KA22 (Контр\_KA22)  & \centering SGF22 & \centering Не предусмотрено \\ Предусмотрено & \centering -- & \centering \arraybackslash -- \\
\hline
\centering Контроль KA23 (Контр\_KA23)  & \centering SGF23 & \centering Не предусмотрено \\ Предусмотрено & \centering -- & \centering \arraybackslash -- \\
\hline
\centering Контроль KA24 (Контр\_KA24)  & \centering SGF24 & \centering Не предусмотрено \\ Предусмотрено & \centering -- & \centering \arraybackslash -- \\
\hline
\centering Контроль KA25 (Контр\_KA25)  & \centering SGF25 & \centering Не предусмотрено \\ Предусмотрено & \centering -- & \centering \arraybackslash -- \\
\hline
\centering Контроль KA26 (Контр\_KA26)  & \centering SGF26 & \centering Не предусмотрено \\ Предусмотрено & \centering -- & \centering \arraybackslash -- \\
\hline
\centering Контроль KA27 (Контр\_KA27)  & \centering SGF27 & \centering Не предусмотрено \\ Предусмотрено & \centering -- & \centering \arraybackslash -- \\
\hline
\centering Контроль KA28 (Контр\_KA28)  & \centering SGF28 & \centering Не предусмотрено \\ Предусмотрено & \centering -- & \centering \arraybackslash -- \\
\hline
\centering Контроль KA29 (Контр\_KA29)  & \centering SGF29 & \centering Не предусмотрено \\ Предусмотрено & \centering -- & \centering \arraybackslash -- \\
\hline
\centering Контроль KA30 (Контр\_KA30)  & \centering SGF30 & \centering Не предусмотрено \\ Предусмотрено & \centering -- & \centering \arraybackslash -- \\
\hline
\centering Контроль KA31 (Контр\_KA31)  & \centering SGF31 & \centering Не предусмотрено \\ Предусмотрено & \centering -- & \centering \arraybackslash -- \\
\hline
\centering Контроль KA32 (Контр\_KA32)  & \centering SGF32 & \centering Не предусмотрено \\ Предусмотрено & \centering -- & \centering \arraybackslash -- \\
\hline
\multicolumn{5}{|c|}{ КП } \\ \hline 
\centering Ввод функции в работу (Ввод\_функции)  & \centering SGF1 & \centering Не предусмотрено \\ Предусмотрено & \centering -- & \centering \arraybackslash -- \\
\hline
\centering Ввод ОБР КА в работу (ОБР\_КА)  & \centering SG1 & \centering Не предусмотрено \\ Предусмотрено & \centering -- & \centering \arraybackslash -- \\
%===t1
\end{longtable}
\normalsize

\end{enumerate}
%%%%%%%%%%%%%%%%%%%%%%%%%%%%%%%%%%%%%%%%%%%%%%%%%%%%%%%%%%%%%%%%%%%%%%  КОНЕЦ ОБЩИЕ СВЕДЕНИЯ %%%%%%%%%%%%%%%%%%%%%%%%%%%%%%%%%%%%%%%%%%%%%%%%%%

%%%%%%%%%%%%%%%%%%%%%%%%%%%%%%%%%%%%%%%%%%%%%%%%%%%%%%%%%%%%%%%%%%%%%%  УПРАВЛЕНИЕ ЗАЗЕМЛЯЮЩИМ НОЖОМ %%%%%%%%%%%%%%%%%%%%%%%%%%%%%%%%%%%%%%%%%%%%%%%%%%
\item Управление заземляющим ножом (УЗН)

\begin{enumerate}[label=\arabic{section}.\arabic{subsection}.\arabic{enumi}.\arabic*, labelsep=4pt, leftmargin=0em, itemindent=65pt, parsep=0pt]

\item 
В функциональном блоке «КП» логика управления заземляющим ножом представлена функцией «УЗН». Функция формирует команды на отключение и включение заземляющего ножа при оперативном управлении.

\item
Команда оперативного отключения «УЗН: Отключить» формируется в следующих случаях:
\begin{itemize}
	\item при появлении команды на оперативное отключение заземляющего ножа («Откл. ЗН ИЧМ»), выбранным управлением по месту (сигнал «Местное» находится в активном состоянии) и наличием разрешающего сигнала от схемы оперативной блокировки «ОБР ЗН: Разрешить управ.»;
	\item при появлении команды на оперативное отключение заземляющего ножа («Откл. ЗН АСУ»), выбранным дистанционным управлением (сигнал «Дистанционное» находится в активном состоянии) и наличием разрешающего сигнала от схемы оперативной блокировки «ОБР ЗН: Разрешить управ.».
\end{itemize}

\item
Прохождение команды на отключение блокируется при наличии следующих сигналов:
\begin{itemize}
	\item отключенное положение заземляющего ножа «ЗН отключен (б/к)» (с контролем сигнала «ЗН включен~(б/к)»);
	\item превышение времени переключения заземляющего ножа «УЗН: Прев. времени~пер.».
\end{itemize}

\item
Команда оперативного включения «УЗН: Включить» формируется в следующих случаях:
\begin{itemize}
	\item при появлении команды на оперативное включение выключателя («Вкл. ЗН ИЧМ»), выбранным управлением по месту (сигнал «Местное» находится в активном состоянии) и наличием разрешающего сигнала от схемы оперативной блокировки «ОБР ЗН: Разрешить управ.»;
	\item при появлении команды на оперативное включение выключателя («Вкл. ЗН АСУ»), выбранным дистанционным управлением (сигнал «Дистанционное» находится в активном состоянии) и наличием разрешающего сигнала от схемы оперативной блокировки «ОБР ЗН: Разрешить управ.».
\end{itemize}

\item
Прохождение команды на включение блокируется при наличии следующих сигналов:
\begin{itemize}
	\item включенное положение заземляющего ножа «ЗН~включен (б/к)» (с контролем сигнала «ЗН~отключен~(б/к)»);
	\item превышение времени переключения заземляющего ножа «УЗН: Прев. времени~пер.».
\end{itemize}

\item
При наличии сигналов «УЗН: Включить» и «УЗН: Отключить» длительностью больше, чем время уставки «T1» формируется сигнал превышения допустимого времени переключения -- «УЗН: Прев. времени пер.». Активное состояние операции включения или отключения сигнализируется выходным сигналом «УЗН: Идет переключение».

\item
Логика функции формирует сигналы состояния выключателя «УЗН: Отключено», «УЗН: Включено», «УЗН: Не определено» и «УЗН: Неисправность», при этом сигналы «УЗН: Неисправность» и «УЗН: Не определено» выдаются с выдержкой времени, задаваемой уставкой Т2 «Тблк».

\item
Ввод функции в работу на этапе параметрирования устройства осуществляется программным переключателем SGF1 «Ввод\_функции». Информация о введенной функции отображается выходным сигналом «УЗН: Ввод».

\item
Алгоритм работы функции выполнен в соответствии с рисунком \ref{kp:img_uzn}. 

\vspace{3mm}
\begin{figure}[H]
\centering
\includegraphics[width=1\textwidth,height=1\textheight,keepaspectratio]{\fbpath/214. КП ТН/_latex/img/uzn.pdf}
\captionof{figure}{Функциональная схема логики УЗН}\label{kp:img_uzn}
\end{figure}


\end{enumerate}

%%%%%%%%%%%%%%%%%%%%%%%%%%%%%%%%%%%%%%%%%%%%%%%%%%%%%%%%%%%%%%%%%%%%%% КОНЕЦ УПРАВЛЕНИЕ ЗАЗЕМЛЯЮЩИМ НОЖОМ %%%%%%%%%%%%%%%%%%%%%%%%%%%%%%%%%%%%%%%%%%%%%%%%%%

%%%%%%%%%%%%%%%%%%%%%%%%%%%%%%%%%%%%%%%%%%%%%%%%%%%%%%%%%%%%%%%%%%%%%%  УПРАВЛЕНИЕ ВЫКАТНЫМ ЭЛЕМЕНТОМ %%%%%%%%%%%%%%%%%%%%%%%%%%%%%%%%%%%%%%%%%%%%%%%%%%

\item Управление выкатным элементом (УВЭ)

\begin{enumerate}[label=\arabic{section}.\arabic{subsection}.\arabic{enumi}.\arabic*, labelsep=4pt, leftmargin=0em, itemindent=65pt, parsep=0pt]

\item 
В функциональном блоке «КП» логика управления выкатным элементом представлена функцией «УВЭ». Функция формирует команды вкатить и выкатить выкатной элемент ячейки фидера при оперативном управлении.

\item
Команда операции выкатить ВЭ «УВЭ: Выкатить» формируется в следующих случаях:
\begin{itemize}
	\item при появлении команды выкатить от местного управления («Выкат. ВЭ ИЧМ»), выбранным управлением по месту (сигнал «Местное» находится в активном состоянии) и наличием разрешающего сигнала от схемы оперативной блокировки «ОБР ВЭ: Разрешить управ.»;
	\item при появлении команды выкатить от дистанционного управления («Выкат. ВЭ АСУ»), выбранным дистанционным управлением (сигнал «Дистанционное» находится в активном состоянии) и наличием разрешающего сигнала от схемы оперативной блокировки «ОБР ВЭ: Разрешить управ.».
\end{itemize}

\item
Прохождение команды выкатить ВЭ блокируется при наличии следующих сигналов:
\begin{itemize}
	\item контрольное положение ВЭ;
	\item превышение времени операции вкатить/выкатить ВЭ «УВЭ: Прев. времени~пер.».
\end{itemize}

\item
Команда операции вкатить ВЭ «УВЭ: Вкатить» формируется в следующих случаях:
\begin{itemize}
	\item при появлении команды вкатить от местного управления («Вкат. ВЭ ИЧМ»), выбранным управлением по месту (сигнал «Местное» находится в активном состоянии) и наличием разрешающего сигнала от схемы оперативной блокировки «ОБР ВЭ: Разрешить управ.»;
	\item при появлении команды вкатить от дистанционного управления («Вкат. ВЭ АСУ»), выбранным дистанционным управлением (сигнал «Дистанционное» находится в активном состоянии) и наличием разрешающего сигнала от схемы оперативной блокировки «ОБР ВЭ: Разрешить управ.».
\end{itemize}

\item
Прохождение команды вкатить ВЭ блокируется при наличии следующих сигналов:
\begin{itemize}
	\item рабочее положение ВЭ;
	\item превышение времени операции вкатить/выкатить ВЭ «УВЭ: Прев. времени~пер.».
\end{itemize}

\item
При наличии сигналов «УВЭ: Вкатить» и «УВЭ: Выкатить» длительностью больше, чем время уставки «T1» формируется сигнал превышения допустимого времени операции -- «УВЭ: Прев. времени пер.». Активное состояние операции включения или отключения сигнализируется выходным сигналом «УВЭ: Идет переключение».

\item
Логика функции формирует сигналы состояния выключателя «УВЭ: Выкачен», «УВЭ: Вкачен», «УВЭ: Не определено» и «УВЭ: Неисправность», при этом сигналы «УВЭ: Неисправность» и «УВЭ: Не определено» выдаются с выдержкой времени, задаваемой уставкой Т2 «Тблк».

\item
Ввод функции в работу на этапе параметрирования устройства осуществляется программным переключателем SGF1 «Ввод\_функции». Информация о введенной функции отображается выходным сигналом «УВЭ: Ввод».

\item
Алгоритм работы функции выполнен в соответствии с рисунком \ref{kp:img_uve}. 

\vspace{3mm}
\begin{figure}[H]
\centering
\includegraphics[width=1\textwidth,height=1\textheight,keepaspectratio]{\fbpath/214. КП ТН/_latex/img/uve.pdf}
\captionof{figure}{Функциональная схема логики УВЭ}\label{kp:img_uve}
\end{figure}


\end{enumerate}
%%%%%%%%%%%%%%%%%%%%%%%%%%%%%%%%%%%%%%%%%%%%%%%%%%%%%%%%%%%%%%%%%%%%%% КОНЕЦ УПРАВЛЕНИЕ ВЫКАТНЫМ ЭЛЕМЕНТОМ %%%%%%%%%%%%%%%%%%%%%%%%%%%%%%%%%%%%%%%%%%%%%%%%%%

%%%%%%%%%%%%%%%%%%%%%%%%%%%%%%%%%%%%%%%%%%%%%%%%%%%%%%%%%%%%%%%%%%%%%%  ОПЕРАТИВНАЯ БЛОКИРОВКА ОБЩЕЕ %%%%%%%%%%%%%%%%%%%%%%%%%%%%%%%%%%%%%%%%%%%%%%%%%%
\item Оперативная блокировка КА (ОБР КА)

\begin{enumerate}[label=\arabic{section}.\arabic{subsection}.\arabic{enumi}.\arabic*, labelsep=4pt, leftmargin=0em, itemindent=65pt, parsep=0pt]

\item
В устройстве реализованы алгоритмы, которые позволяет контролировать положение КА и выдавать сигналы на разрешение управления выкатным элементом (ВЭ) и заземляющим ножом (ЗН).

\item
Поясняющая схема первичной сети для алгоритмов оперативной блокировки КА приведена на рисунке \ref{kp:img_obrka}.

\vspace{1mm}
\begin{figure}[H]
\centering
\includegraphics[width=0.65\textwidth,height=0.65\textheight,keepaspectratio]{\fbpath/214. КП ТН/_latex/img/obrka.pdf}
\captionof{figure}{Поясняющая схема первичной сети для алгоритмов оперативной блокировки КА}\label{kp:img_obrka}
\end{figure}

\end{enumerate}
%%%%%%%%%%%%%%%%%%%%%%%%%%%%%%%%%%%%%%%%%%%%%%%%%%%%%%%%%%%%%%%%%%%%%% КОНЕЦ ОПЕРАТИВНАЯ БЛОКИРОВКА ОБЩЕЕ %%%%%%%%%%%%%%%%%%%%%%%%%%%%%%%%%%%%%%%%%%%%%%%%%%

%%%%%%%%%%%%%%%%%%%%%%%%%%%%%%%%%%%%%%%%%%%%%%%%%%%%%%%%%%%%%%%%%  Оперативная блокировка заземляющего ножа (ОБР ЗН) %%%%%%%%%%%%%%%%%%%%%%%%%%%%%%%%%%%%%%%%%%%%%%%%%%
\item Оперативная блокировка заземляющего ножа (ОБР ЗН)

\begin{enumerate}[label=\arabic{section}.\arabic{subsection}.\arabic{enumi}.\arabic*, labelsep=4pt, leftmargin=0em, itemindent=65pt, parsep=0pt]

\item
В функциональном блоке «КП» логика оперативной блокировки заземляющего ножа представлена функцией «ОБР ЗН». Условие ввода в работу разрешения управления ЗН является наличие входного сигнала «УЗН: Ввод». Контроль условий блокировки управления ВЭ осуществляется с помощью уставки SG1. Если уставка SG1 (см. таблицу \ref{kp:tbl_gen}) выставлена в положение «Не предусмотрено», то управление ЗН осуществляется без учета всех блокировок.

\item
Принудительное деблокирование управления заземляющим ножом от оперативной блокировки осуществляется активацией входного сигнала от ключа «Программ. деблок. ЗН».

\item
Разрешение управления ЗН формируется при выполнении необходимых условий:
\begin{itemize}
	\item отсутствие процесса переключения коммутационных аппаратов (сигнал «КП: Идет переключение»);
	\item активное состояние сигнала «ЗН: Включен» для разрешения команды отключения ЗН или одновременное наличие следующих сигналов для разрешения команды включения ЗН: «ВЭ: Выкачен», «ЗН: Отключен», «КА3 Отключен», «КА4 Отключен», если введен контроль данных сигналов соответствующими уставками.
\end{itemize}

\item\sloppy
Внешние дискретные входы назначаются в проекте непосредственно под схему объекта. Блокировка управления ЗН выполняется при несоблюдении какого-либо из условий, предусмотренных алгоритмом функции.

\item
Алгоритм работы функции выполнен в соответствии с рисунком \ref{kp:img_obrzn}. 

\vspace{1mm}
\begin{figure}[H]
\centering
\includegraphics[width=1\textwidth,height=1\textheight,keepaspectratio]{\fbpath/214. КП ТН/_latex/img/obrzn.pdf}
\captionof{figure}{Функциональная схема логики ОБР ЗН}\label{kp:img_obrzn}
\end{figure}


\end{enumerate}
%%%%%%%%%%%%%%%%%%%%%%%%%%%%%%%%%%%%%%%%%%%%%%%%%%%%%%% КОНЕЦ Оперативная блокировка заземляющего ножа (ОБР ЗН) %%%%%%%%%%%%%%%%%%%%%%%%%%%%%%%%%%%%%%%%%%%%%%%%

%%%%%%%%%%%%%%%%%%%%%%%%%%%%%%%%%%%%%%%%%%%%%%%%%%%%%%%%%  Оперативная блокировка выкатного элемента (ОБР ВЭ) %%%%%%%%%%%%%%%%%%%%%%%%%%%%%%%%%%%%%%%%%%%%%%%%%%
\item Оперативная блокировка выкатного элемента (ОБР ВЭ)

\begin{enumerate}[label=\arabic{section}.\arabic{subsection}.\arabic{enumi}.\arabic*, labelsep=4pt, leftmargin=0em, itemindent=65pt, parsep=0pt]

\item
В функциональном блоке «КП» логика оперативной блокировки выкатного элемента представлена функцией «ОБР ВЭ». Условие ввода в работу разрешения управления ВЭ является наличие входного сигнала «УВЭ: Ввод». Контроль условий блокировки управления ВЭ осуществляется с помощью уставки SG1. Если уставка SG1 (см. таблицу \ref{kp:tbl_gen}) выставлена в положение «Не предусмотрено», то управление ВЭ осуществляется без учета всех блокировок.

\item
Принудительное деблокирование управления выкатным элементом от оперативной блокировки осуществляется активацией входного сигнала от ключа «Программ. деблок. ВЭ».

\item
Разрешение управления ВЭ формируется при выполнении необходимых условий:
\begin{itemize}
	\item отсутствие процесса переключения коммутационных аппаратов (сигнал «КП: Идет переключение»);
	\item активное состояние сигнала «ВЭ: Вкачен» для разрешения команды выкатить ВЭ или одновременное наличие следующих сигналов для разрешения команды вкатить ВЭ: «ВЭ: Выкачен», «ЗН: Отключен», «КА1 Отключен», «КА2 Отключен», если введен контроль данных сигналов соответствующими уставками.
\end{itemize}

\item\sloppy
Внешние дискретные входы назначаются в проекте непосредственно под схему объекта. Блокировка управления ВЭ выполняется при несоблюдении какого-либо из условий, предусмотренных алгоритмом функции.

\item
Алгоритм работы функции выполнен в соответствии с рисунком \ref{kp:img_obrve}.

\vspace{1mm}
\begin{figure}[H]
\centering
\includegraphics[width=1\textwidth,height=1\textheight,keepaspectratio]{\fbpath/214. КП ТН/_latex/img/obrve.pdf}
\captionof{figure}{Функциональная схема логики ОБР ВЭ}\label{kp:img_obrve}
\end{figure}

\end{enumerate}
%%%%%%%%%%%%%%%%%%%%%%%%%%%%%%%%%%%%%%%%%%%%%%%%%%%%%%%%%%%%%% КОНЕЦ Оперативная блокировка выкатного элемента (ОБР ВЭ) %%%%%%%%%%%%%%%%%%%%%%%%%%%%%%%%%%%%%%%%%%%%%%%%%%


\end{enumerate}
